\chapter{Digital signal processing}
\label{ch:digital_signal_proc}
Electrochemical impedance spectra are obtained by processing time-domain voltage and current signals measured under controlled perturbations.
As a consequence, impedance is not directly observed, but estimated through signal processing operations that transform finite, noisy, and discretely sampled data into a frequency-domain representation.

The accuracy, resolution, and physical validity of impedance measurements therefore depend not only on the electrochemical system and the applied excitation, but also on the signals processing.
While system theory defines impedance as the frequency response of a linear system, digital signal processing determines how this response can be experimentally accessed in practice.

This chapter introduces the foundamental signal processing concepts required to manipulate discrete digital signals, including the discrete Fourier transform, spectral leakage, and noise considerations.
These tools form the basis for understanding both classical impedance spectroscopy and the advanced time-resolved methods discussed in the following chapters.

\section{Sampling and quantization}
In the real world, any quantity changes continuously in time. 
When performing an experiment, the scientist measures the values of those quantities with instrumentation at specific time intervals.
In practice, the electronic circuit that digitally records analog signals is called an analog-to-digital converter (ADC) and is a staple in digital signal processing.

Any digital measurement system samples a continuous signal $x(t)$ at discrete times $t=nT_s$, producing the sampled discrete signal $x[n]$; where $T_s$ is the sampling time. 
The most important parameter to consider in sampling is the time interval or its reciprocal: the sampling frequency $f_s = 1/T_s$.
To analyze the spectral content of a signal, the sampling frequency must be at least twice the highest oscillation frequency. This is known as the Nyquist--Shannon sampling theorem. Sampling below this threshold causes higher-frequency components to "fold back" into lower frequencies—a phenomenon known as aliasing—which corrupts the true spectral content of the signal.

\section{Discrete Fourier transform}
After acquiring a digital signal, the role of digital signal processing is to analyze or act on the frequency components of the signal.
As introduced in the previous chapter, the Fourier transform is used for spectral analysis of signals, which is the process of decomposing a time-domain signal into its constituent frequencies and quantifying the amplitude and phase associated with each.

One converts a signal $x(t)$ or $x[n]$ into $X(\omega)$ or $X[k]$, revealing how its energy or variance is distributed across frequency. 

Since digital signals are discrete vectors, the discrete version of the Fourier transform is used, called the discrete Fourier transform (DFT):
\begin{equation}
X[k]=\sum_{n=0}^{N-1} x[n] e^{-j2\pi kn/N}.
\end{equation}
The inverse discrete Fourier transform (iDFT) takes the form:
\begin{equation}
x[n]= \frac{1}{N} \sum_{k=0}^{N-1} X[k] e^{j2\pi kn/N}.
\end{equation}

The discrete Fourier transform of a signal is calculated in a computer using the fast Fourier transform (FFT) algorithm, the most efficient implementation developed by Cooley and Tukey.
The computational complexity is $O(N\log N)$ compared to the $O(N^2)$ of direct summation, although it requires the signal length to be a power of 2, i.e., $N = 2^n$ for some integer $n$.

% \subsection{Short-time Fourier Transform (STFT)}

% [Put some introduction here on why the STFFT is important in genreal and why for this work.]

% To analyse drifting or time-varying systems one applies the STFT:
% \[
% X(m,k)=\sum_{n=0}^{N-1} x[n+mR]\, w[n]\, e^{-j2\pi kn/N},
% \]
% where $m$ indexes time windows and $w[n]$ is a window function.  
% The STFT reveals how the spectral content evolves in time and provides the foundation for non-stationary impedance estimation discussed in later chapters.

% In real world applications of signal processing, the spectral information changes with time. A few examples are the real-time processing of music (digital equalization), human voice (speech recognition) or images (object recognition). To analyze this signals, the approach is to perform the Fourier Trasnform on a small part of the signal at the time, with a process called windowing. The method is called short-time Fourier Trasnform (STFT) and extends the classic FT to non-stationary signals by sliding a finite window w[n] along the data and computing a local DFT on each segment. Formally, for a discrete signal x[n] and window shift R:

% $$
% X(m,k) = \sum _{n=0} ^{n-1} x[n + mR] \cdot w[n] \cdot e^{-j2\pi kn/N},
% $$

% where m indexes time frames and k frequency bins. While the short-time Fourier Trasnfomr can effectively capture the spectral contect of transient processes, the accuracy of the result depends on careful design of the calculation. First, there is a time-frequency resolution trade-off, known also as Gabor uncenrtainty principle $\Delta t \cdot \Delta f $ where $\Delta t$ is the windows lenght and $\Delta f$ is the inverse of it. Equality is reached only by a Gaussian window, which is therefore “optimal” in the sense of joint time-frequency localization. The window length N controls the balance: a short window localizes well in time (small $\Delta t$) but its Fourier transform is broad (large $\Delta f$), so fast transients are captured but spectral peaks smear, while a long window gives fine frequency discrimination (small $\Delta f$) at the expense of temporal precision (large $\Delta t$), so you resolve closely spaced frequencies but miss rapid changes. Another important aspect is the window shape and the overlap, the latter derived by the windows shift value $R$ in the STFT definition.  To produce smooth progression between following windowing, an overlap is used of 50\% or 75\%. 

% These concepts are essential for the discussion on the methods to estimate the non-stationary impedance given later in the text and they must be carefully taken into consideration depending on the characteristic frequency of evolution of the given signal.


\subsection{Spectral leakage}

Spectral leakage is the effect in the Discrete Fourier Transform of spectral components "leaking" or "spilling" into neighboring frequency bins, producing spurious components where the spectrum baseline should be flat (at noise level).
Multiple causes produce spectral leakage.
The most common cause is a discrete signal $x[n]$ that does not contain an integer number of periods of the oscillatory components.
Other causes of spectral leakage derive from practical hardware limitations or user errors in the experimental design, such as:
\begin{itemize}
    \item lack of synchronization between ADC and DAC clocks;
    \item instability of the ADC (data acquisition) or DAC (signal generator) clock;
    \item ADC jitter and slew-rate;
    \item low resolution of the ADC compared to the oscillation amplitude, which may produce spurious harmonics;
    \item aliasing due to improper choice of sampling frequency without adequate filtering.
\end{itemize}

\subsection{Noise}

Noise denotes random, undesired fluctuations superimposed on the signal originating from physical sources, instrumentation (quantization, clock jitter) or environment (EM interference). 

Signal-to-noise ratio (SNR) quantifies the relative power between the signal components and the noise:
$$
SNR = \frac{P_{\mathrm{signal}}}{P_{\mathrm{noise}}}
$$
and dictates how reliably the signal can be detected or processed. 

When increasing the excitation amplitude is not possible, averaging is an alternative tool for improving SNR: by repeating the same measurement or subdividing a long record into M independent segments, uncorrelated noise averages toward zero while the coherent signal adds constructively.
This concept is used in electrochemical instrumentation where impedance spectroscopy is implemented.
The user can choose the number of sinusoidal wave periods to average.
The net SNR gain scales with the square root of the number of independent measurements.
In the time domain, this is called ensemble or block averaging; in the frequency domain, it appears as spectral averaging.
Averaging can alternatively be obtained by increasing the acquisition time of the periodic signal or the window length.
In fact, narrowing the spectral bin bandwidth reduces the noise power in it linearly (i.e., doubling the signal length halves $\Delta f$, and thus the power in each bin, while the coherent signal amplitude in the bin does not change).

\section{Filters}

A filter in signal processing is a system that selectively modifies an input signal's amplitude and/or phase as a function of frequency and it is a core subject of signal processing.

By their effect on frequency content, filters are classified as low-pass (passing frequencies below a cutoff $f_c$), high-pass (passing frequencies above $f_c$), band-pass (passing a band between two cutoffs), or band-stop/notch (rejecting a specific band).  
  
Filters can be distinguished between analog, a physical circuit in the signal acquisition pipeline, or digital, applied on the digitized signal by a micro-controller or computer.

There are different types of digital implementation of filters, but for the purpose of this work it is only necessary to consider the effect on the frequency content of a signal and performance metrics such as passband ripples (amplitude variation within the passband), stopband attenuation (suppression of unwanted bands), and transition bandwidth (width of the roll-off region).  

% \subsection{Filters in DSP (to be corrected)}
% Although system identification in electrochemsitry for the frequency domain employs cyclic perturbation (i.e. single sine waves or multi-sines), it will be usefull in the following discussio in this text to introduce some aspect of window functions. 
% The most known are Hanning, Hamming and Blackman  functions characterised by three main properties:
% \begin{itemize}
%     \item roll-off and characteristic frequency
%     \item passband flatness
%     \item group delay
% \end{itemize}
% The main characteristic is called roll-off which means "how fast" the amplitude of the original signal is attenuated. From the view of the frequency domain, the part of the filter that attenuates the signal are colled sidelobes. The characteristic frequency, also called cut-off, is the frequency at which the attenuation is $\sqrt{2}$ (\~= -3 dB) of the maximum value. Different windows function have steeper or smoother roll-off and this is important when realizing, for example, a low pass filter using the same function. Passband flatness is how uniform is the magnitude response across the main lobe (the passband) to avoid any distortion of the amplitude.
% It is also important, especially for the complex signals like the impedance, to evaluate also the phase shift or group delay that must be flat in the frequency region of analysis. Windows have usually a zero or very low attenuation in the middle that translates to a very well analysis of the spectral characteristics of the middle point of the signal. This is somehow translated to an "instantaneous information" even if the windows has finite length.
% Usually it is impossible to have optimal values for this parameters but a compromise has to be made.

% \section{Removal of trends}
% Baseline or trend removal, often called detrending, is the process of eliminating slow-varying offsets or drifts from a signal so that its remaining fluctuations can be treated as (approximately) zero-mean and stationary.  
% In electrochemical measurements, trends arise from the drift the system is undertaking, its thermodynamic state and relaxation processes.
% If left uncorrected they bias spectral estimates, manifesting as excessive low-frequency power or leakage and undermines the excitation frequency analysis.  
% Common detrending techniques include subtracting the signals long-term mean or a low-order polynomial fit, applying a high-pass filter with a cutoff below the band of interest (using a linear-phase FIR to avoid phase distortion), or employing moving-average/median filters to track and remove slow baseline wander. 
% In all cases, one must choose parameters (filter cutoff, polynomial order, window length) so that genuine signal dynamics are preserved while the unwanted trend is effectively suppressed.

% In the analysis of signals trough Fourier Analysis, the removal of trends mitigate an effect called Gibbs phenomenon that is the presence of an infinite series of harmonics produce by sharp transition in the signals caused by its dynamic (drift) as well as instrumental coming from the digitization process (an artifact known as ringing).

\section{Signal synthesis}

To perform certain experiments, one might want to perturb the system with non-constant signals such as pulses, ramps, sinusoids, or other arbitrary shapes.
This requires designing a digital waveform and generating it in analog form using electronic instrumentation.
This process is called digital-to-analog conversion, and the electronic component performing it is the digital-to-analog converter (DAC).
Several techniques exist for implementing a digital-to-analog converter on hardware with different complexity and properties; the simplest is \emph{direct digital synthesis} (DDS).
This method relies on a 'look-up table' in which the desired waveform is stored as a table of value-time pairs and at constant time intervals the value is read and converted to an analog voltage.
In arbitrary waveform generators, the table of points is usually stored in limited buffer memory of a few thousand samples.
The generated signal is typically passed through a series of filters to smooth it and avoid producing harmonics or causing spectral leakage.
The synthesis procedure is crucial for the quality of the excitation in system identification.

The length of the arbitrary waveform can be extended by storing the signal in larger memory and sending it directly to the digital-to-analog converter using the direct memory access (DMA) protocol; however, this approach requires dedicated micro-controller units and knowledge of digital signal processor implementation that may be outside the typical background of electrochemists.

\section{Conclusion}
The mathematical framework introduced in this chapter enables the transformation and analysis of time-domain signals in the frequency domain.
However, the application of these tools to electrochemical impedance spectroscopy requires additional assumptions regarding system linearity, time invariance, and experimental excitation.
In particular, when the electrochemical system evolves on time scales comparable to the measurement duration, the direct application of classical frequency-domain techniques becomes inadequate.
The following chapter therefore introduces the methodological extensions required to estimate impedance under non-stationary conditions.